\documentclass[a4paper,12pt]{article}

\usepackage{alltt, fancyvrb, url}
\usepackage{graphicx}
\usepackage[utf8]{inputenc}
\usepackage{float}
\usepackage{hyperref}
\usepackage[italian]{babel}
\usepackage[table]{xcolor}
\usepackage{array}
\usepackage{multirow}
\usepackage{titlesec}
\usepackage{listings}
\usepackage[italian]{cleveref}
\usepackage{acronym}
\usepackage{subcaption}

\renewcommand{\thesection}{\arabic{section}}

\sloppy
\setlength{\parindent}{0pt}

\title{\textbf{Responsabilità dei task}}
\author{}
\date{}

\begin{document}
\maketitle

Il team di sviluppo è composto da 10 membri, di cui 5 sono specializzati nel frontend e 5 nel backend, e i team sono
organizzati in maniera orizzontale. Entrambi i gruppi sono composti da sviluppatori senior e junior, per questo motivo
i responsabili delle attività saranno gli sviluppatori senior, del relativo ambito.
Per ogni attività si realizza una matrice di assegnazione delle responsabilità (seguendo lo schema \textit{RASCI}).
Dunque per ogni task si assegna un responsabile ed eventuali altri ruoli (Accountable, Support, Consulted, Informed).

\end{document}